
\subsection{DSP Processing Pipeline}

Pipeline B applies two sequential DSP stages before feature extraction. The signal flow is:

\begin{center}
\begin{tabular}{ccccccc}
Preprocessed & $\xrightarrow{\text{Stage 1}}$ & FIR Bandpass & $\xrightarrow{\text{Stage 2}}$ & Pre-emphasis & $\xrightarrow{}$ & MFCC \\
audio & & 300--3400 Hz & & $\alpha = 0.97$ & & extraction \\
\end{tabular}
\end{center}

\noindent \textbf{Stage 1 (FIR Filter)} isolates the speech-relevant frequency band, removing low-frequency hum ($<300$ Hz) and high-frequency noise ($>3400$ Hz). \textbf{Stage 2 (Pre-emphasis)} compensates for the natural spectral tilt of speech, boosting high-frequency components that carry consonant and speaker-discriminative information. Together, these two stages increase the Signal-to-Noise Ratio (SNR) and flatten the spectrum before MFCC computation.

\subsection{FIR Design}

This section presents the complete set of formulas for FIR filter design using the window method (as implemented in \texttt{firwin}).

\subsubsection{Ideal Impulse Response ($h_{\text{ideal}}$)}
The ideal bandpass filter is constructed by subtracting the impulse response of a lower-frequency low-pass filter from a higher-frequency low-pass filter:
\begin{equation}
h_{\text{ideal}}[n] = \frac{\sin\left(\omega_{\text{high}} \left(n - \frac{M}{2}\right)\right)}{\pi \left(n - \frac{M}{2}\right)} - \frac{\sin\left(\omega_{\text{low}} \left(n - \frac{M}{2}\right)\right)}{\pi \left(n - \frac{M}{2}\right)}
\end{equation}
Where:
\begin{itemize}
  \item $M$: The filter order ($M = \text{numtaps} - 1$).
  \item $n$: The sample index, ranging from $0$ to $M$.
  \item $\omega_{\text{low}}, \omega_{\text{high}}$: normalized angular frequencies (rad/sample), calculated as
  \begin{equation}
  \omega = \frac{2\pi f}{f_s}
  \end{equation}
  \item $\left(n - \frac{M}{2}\right)$: the time shift that ensures the filter is symmetric and has a linear phase.
\end{itemize}


\subsubsection{Window Function ($w[n]$)}
To truncate the infinite ideal impulse response into a finite one, we multiply it by a Hamming window:
\begin{equation}
w[n] = 0.54 - 0.46 \cos\left(\frac{2\pi n}{M}\right)
\end{equation}

\subsubsection{Final Filter Coefficients ($h[n]$)}
These are the filter taps returned by \texttt{scipy.signal.firwin} in the coefficient array:
\begin{equation}
h[n] = h_{\text{ideal}}[n] \cdot w[n]
\end{equation}

\subsubsection{Handling the Singular Point}
At the center of the filter, where $n = \frac{M}{2}$, the denominator becomes zero. By L'H\^{o}pital's rule, the value at this point is defined as:
\begin{equation}
h_{\text{ideal}}\left[\frac{M}{2}\right] = \frac{\omega_{\text{high}} - \omega_{\text{low}}}{\pi}
\end{equation}

\subsection{Pre-emphasis (High-Frequency Boosting)}

\textbf{The Problem (Spectral Tilt):} In natural speech, the energy of the signal decreases as the frequency increases (typically at a rate of $-6$ dB/octave). Consequently, high-frequency components (such as consonants and sibilants) are significantly weaker than low-frequency components (vowels). 

\textbf{The Solution:} To balance the power spectrum, we apply a first-order pre-emphasis high-pass filter to the 1-D audio signal. This process amplifies the high-frequency content, making the signal more uniform across the frequency domain. 

The mathematical representation of this filter in the discrete-time domain is:
\begin{equation}
    y'[n] = y[n] - \alpha \cdot y[n-1]
\end{equation}

\textbf{Parameters:}
\begin{itemize}
    \item $y$: 1-D array representing the input audio signal.
    \item $\alpha$: Pre-emphasis coefficient in $[0, 1)$. Typical values range from $0.95$ to $0.97$ (most commonly $0.97$).
\end{itemize}

\textbf{Mechanism \& Impact:} 
By subtracting a large percentage of the previous sample from the current one, the filter suppresses slowly changing components (low frequencies) while highlighting rapid changes (high frequencies). This results in \textit{spectral flattening}, ensuring that the subsequent Mel-filterbank and DCT steps can efficiently extract meaningful features from the entire frequency range, rather than being dominated by low-frequency energy.